
\documentclass[10pt]{article}
\usepackage[a4paper, tmargin=0.75in, lmargin=0.80in, rmargin=0.80in, bmargin=1in]{geometry}
\usepackage{hyperref}
\usepackage{dsfont}
\usepackage{amsfonts}
\usepackage{ragged2e}
\usepackage{mathtools}
\usepackage{enumitem}   
\usepackage{MnSymbol}%
\usepackage{wasysym}%
\usepackage{relsize}%
\usepackage{xcolor}

\usepackage{amsmath}
\usepackage{upgreek}
\usepackage{mathrsfs}
\usepackage{blindtext}
\usepackage{listings}
\usepackage{xcolor}


%New colors defined below
\definecolor{codegreen}{rgb}{0,0.6,0}
\definecolor{codegray}{rgb}{0.5,0.5,0.5}
\definecolor{codepurple}{rgb}{0.58,0,0.82}
\definecolor{backcolour}{rgb}{0.95,0.95,0.92}

\lstset{
    literate={ĉ}{{\^c}}1
        {Ĉ}{{\^C}}1
}

%Code listing style named "mystyle"
\lstdefinestyle{mystyle}{
  backgroundcolor=\color{backcolour}, commentstyle=\color{codegreen},
  keywordstyle=\color{magenta},
  numberstyle=\tiny\color{codegray},
  stringstyle=\color{codepurple},
  basicstyle=\ttfamily\footnotesize,
  breakatwhitespace=false,         
  breaklines=true,                 
  captionpos=b,                    
  keepspaces=true,                 
  numbers=left,                    
  numbersep=5pt,                  
  showspaces=false,                
  showstringspaces=false,
  showtabs=false,                  
  tabsize=2
}

\lstset{style=mystyle}


\usepackage{natbib}
\usepackage{graphicx}
\pagestyle{empty}
\renewcommand\thesubsection{\Alph{subsection}}
\DeclareMathOperator*{\limop}{\lim}




%%%%%%%%%%%%%%%%%%%%%%%%%%%%%%%%%%%%%%%%%%%%%%%%%%
%%%%%%%%%%%%%%%%%%%%%%%%%%%%%%%%%%%%%%%%%%%%%%%%%%
%%%%%%%%%%%%%%%%%%%%%%%%%%%%%%%%%%%%%%%%%%%%%%%%%%
%%%%%%%%%%%%%%%%%%%%%%%%%%%%%%%%%%%%%%%%%%%%%%%%%%
% ENTER SOME IMPORTANT INFORMATION
%%%%%%%%%%%%%%%%%%%%%%%%%%%%%%%%%%%%%%%%%%%%%%%%%%
%%%%%%%%%%%%%%%%%%%%%%%%%%%%%%%%%%%%%%%%%%%%%%%%%%
%%%%%%%%%%%%%%%%%%%%%%%%%%%%%%%%%%%%%%%%%%%%%%%%%%
%%%%%%%%%%%%%%%%%%%%%%%%%%%%%%%%%%%%%%%%%%%%%%%%%%
\newcommand{\studentname}{Luis Alejandro Rubiano}
\newcommand{\studentnumber}{202013482}
\newcommand{\researchcentre}{la.rubiano@uniandes.edu.co}
\newcommand{\projecttitle}{Convex Optimization}
\newcommand{\supervisor}{Mauricio Junca}
\newcommand{\xRightarroww}[2][]{\ext@arrow 0359\Rightarrowfill@{#1}{#2}}

\newcommand\setItemnumber[1]{\setcounter{enum\romannumeral\@enumdepth}{\numexpr#1-1\relax}}

\newcommand\restr[2]{{% we make the whole thing an ordinary symbol
  \left.\kern-\nulldelimiterspace % automatically resize the bar with \right
  #1 % the function
  \littletaller % pretend it's a little taller at normal size
  \right|_{#2} % this is the delimiter
  }}

\newcommand{\littletaller}{\mathchoice{\vphantom{\big|}}{}{}{}}


%%%%%%%%%%%%%%%%%%%%%%%%%%%%%%%%%%%%%%%%%%%%%%%%%%
%%%%%%%%%%%%%%%%%%%%%%%%%%%%%%%%%%%%%%%%%%%%%%%%%%
%%%%%%%%%%%%%%%%%%%%%%%%%%%%%%%%%%%%%%%%%%%%%%%%%%
%%%%%%%%%%%%%%%%%%%%%%%%%%%%%%%%%%%%%%%%%%%%%%%%%%
%%%%%%%%%%%%%%%%%%%%%%%%%%%%%%%%%%%%%%%%%%%%%%%%%%
%%%%%%%%%%%%%%%%%%%%%%%%%%%%%%%%%%%%%%%%%%%%%%%%%%

\begin{document}

\begin{center}
{\Huge{Homework 3, Convex Optimization}} \\
\vspace{2mm}
{\Large{Departamento de Matemáticas, Universidad de los Andes}} \\
\vspace{1mm}
{\Large{Colombia}}
\end{center}

\vspace{5mm}
\hrule
\vspace{1mm}
\hrule

\vspace{3mm}
\begin{tabular}{ll} 
Student's name:           	        & {\studentname}   \\ 
Student's code: 	        & {\studentnumber} \\ 
Email address: 	        & {\researchcentre}  \\ 
Class: 	& {\projecttitle}  \\ 
Instructor: 	    & {\supervisor}  \\ 
\end{tabular}

\vspace{3mm}
\hrule
\vspace{1mm}
\hrule

\hfill \break %SPACE

\section{Let $K_1$ and $K_2$ be convex cones in a vector space $E$. 
Show that $K_1 + K_2$ is a convex cone, $K_1 + K_2 \subseteq \text{co}(K_1 \cup K_2)$, 
if both cones 
contain the origin, then $K_1 + K_2 = \text{co}(K_1 \cup K_2)$. 
Ex. 4.3. Foundations of Optimization, Güler.}

\paragraph{Proof:}

Let $k_{a1} + k_{a2}, k_{b1} + k_{b2} \in K_1 + K_2$ and $t > 0$, then $t (k_{a1} + k_{a2}) = t k_{a1} + t k_{a2}$, since $t k_{a1} \in K_1$ and $t k_{a2} \in K_2$, because $K_1$ and $K_2$ are convex cones. Therefore, $t (k_{a1} + k_{a2}) \in K_1 + K_2$.\\

Also $(k_{a1} + k_{a2}) + (k_{b1} + k_{b2}) = (k_{a1} + k_{b1}) + (k_{a2} + k_{b2})$, since $k_{a1} + k_{b1} \in K_1$ and $k_{a2} + k_{b2} \in K_2$, because $K_1$ and $K_2$ are convex cones. Therefore, $(k_{a1} + k_{a2}) + (k_{b1} + k_{b2}) \in K_1 + K_2$.\\

Thus, $K_1 + K_2$ is a convex cone.\\

Let $k_{1} + k_{2} \in K_1 + K_2$, then $k_{1} + k_{2} = \frac{1}{2} (2 k_{1}) + \frac{1}{2}(2 k_{2})$, since $2 k_{1} \in K_1$ and $2 k_{2} \in K_2$, because $K_1$ and $K_2$ are convex cones. Therefore, since $\frac{1}{2} + \frac{1}{2} = 1$, then $k_{1} + k_{2} \in \text{co}(K_1 \cup K_2)$. i.e. $K_1 + K_2 \subseteq \text{co}(K_1 \cup K_2)$.\\

If both cones contain the origin, then $K_1 = K_1 + 0 \in K_1 + K_2$ and $K_2 = 0 + K_2 \in K_1 + K_2$. Therefore, since $K_1 + K_2$ is convex, and $\text{co}(K_1 \cup K_2)$ is the smallest convex set containing $K_1 \cup K_2$, then $K_1 + K_2 = \text{co}(K_1 \cup K_2). \ \square$

\hfill \break %SPACE

\setcounter{section}{2}

\section{Which of the following sets are convex? Ex. 2.12. Convex Optimization, Boyd.}


\begin{enumerate}[label=(\alph*)]
  \item A \textit{slab}, i.e. a set of the form $\{x \in \mathbb{R}^n \ | \ \alpha \leq a^Tx \leq \beta\}$, where $a \in \mathbb{R}^n$ and $\alpha, \beta \in \mathbb{R}$.
\end{enumerate}

\paragraph{Solution:}


Let $H_1 = \{x \in \mathbb{R}^n \ | \ \alpha \leq a^Tx \}$ and $H_2 = \{x \in \mathbb{R}^n \ | \ a^Tx \leq \beta\}$, then $H_1$ and $H_2$ are halfspaces, and halfspaces are convex. Therefore, the intersection of two convex sets is convex, i.e. $H_1 \cap H_2$ is convex. Thus, a slab is convex. $\square$




\begin{enumerate}[label=(\alph*)]
  \setcounter{enumi}{1}
  \item A \textit{rectangle}, i.e. a set of the form $\{x \in \mathbb{R}^n \ | \ \alpha_i \leq x_i \leq \beta_i, i = 1, \dots, n\}$, where $\alpha_i, \beta_i \in \mathbb{R}$.
\end{enumerate}

\paragraph{Solution:}

Let $H_{i1} = \{x \in \mathbb{R}^n \ | \ \alpha_i \leq x_i \}$ and $H_{i2} = \{x \in \mathbb{R}^n \ | \ x_i \leq \beta_i\}$, then $H_{i1}$ and $H_{i2}$ are halfspaces, and halfspaces are convex. Therefore, the intersection of two convex sets is convex, i.e. $H_{i1} \cap H_{i2}$ is convex. Thus, 

$$\{x \in \mathbb{R}^n \ | \ \alpha_i \leq x_i \leq \beta_i, i = 1, \dots, n\} = \bigcap_{i=1}^n \left(H_{i1} \cap H_{i2}\right)$$ 

is convex because is the finite intersection of convex sets. $\square$


\begin{enumerate}[label=(\alph*)]
  \setcounter{enumi}{2}
  \item A \textit{wedge}, i.e. $\{ x \in \mathbb{R}^n \ | \ a_1^Tx \leq b_1, a_2^Tx \leq b_2\}$, where $a_1, a_2 \in \mathbb{R}^n$ and $b_1, b_2 \in \mathbb{R}$.
\end{enumerate}

\paragraph{Solution:}

Let $H_1 = \{ x \in \mathbb{R}^n \ | \ a_1^Tx \leq b_1\}$ and $H_2 = \{ x \in \mathbb{R}^n \ | \ a_2^Tx \leq b_2\}$, then $H_1$ and $H_2$ are halfspaces, and halfspaces are convex. Therefore, the intersection of two convex sets is convex, i.e. $H_1 \cap H_2$ is convex. Thus, a wedge is convex. $\square$ 




\begin{enumerate}[label=(\alph*)]
  \setcounter{enumi}{3}
  \item The set of points closer to a given point than a given set, i.e. $\{x \ | \ \| x-x_0 \|_2 \leq \| x - y \|_2, \forall y \in S\}$, where $x_0 \in \mathbb{R}^n$ and $S \subseteq \mathbb{R}^n$.
\end{enumerate}

\paragraph{Solution:}

For a fixed $y_0 \in S$, the set $\{x \ | \ \| x-x_0 \|_2 \leq \| x - y_0 \|_2\}$ is a halfspace, thus 

$$ \{x \ | \ \| x-x_0 \|_2 \leq \| x - y \|_2, \forall y \in S\} = \bigcap_{y \in S} \{x \ | \ \| x-x_0 \|_2 \leq \| x - y \|_2\}$$

is the intersection of halfspaces, and the intersection of convex sets is convex. Therefore, the set of points closer to a given point than a given set is convex. $\square$




\begin{enumerate}[label=(\alph*)]
  \setcounter{enumi}{4}
  \item The set of points closer to one set than another, i.e. $\{ x \ | \ \text{dist}(x, S) \leq \text{dist}(x, T)\}$, where $S, T \subseteq \mathbb{R}^n$ and $\text{dist}(x, S) = \inf_{z \in S} \| x - z \|_2$.
\end{enumerate}

\paragraph{Solution:}

This set is not convex in general. For example, let $S = \{0,1\}$ and $T = \{\frac{1}{2}\}$, then the set of points closer to $S$ than $T$ is $\{ x \ | \ \text{dist}(x, S) \leq \text{dist}(x, T)\} = (-\infty, \frac{1}{4}] \cup [\frac{3}{4}, \infty)$, which is not convex, because 
the convex in $\mathbb{R}$ are intervals. $\square$


\begin{enumerate}[label=(\alph*)]
  \setcounter{enumi}{5}
  \item The set $\{ x \ | \ x+S_2 \subseteq S_1\}$, where $S_1, S_2 \subseteq \mathbb{R}^n$ with $S_1$ convex.
\end{enumerate}

\paragraph{Solution:}

This set can be written as
 $$\{ x \ | \ x+S_2 \subseteq S_1\} = \bigcap_{y \in S_2} \{ x \ | \ x+y \in S_1\} = \bigcap_{y \in S_2} (S_1 - y)$$

 $S_1 - y$ is a translation of $S_1$ by $-y$, and translations of convex sets are convex. Therefore, the set $\{ x \ | \ x+S_2 \subseteq S_1\}$ is convex because is the intersection of convex sets. $\square$


\begin{enumerate}[label=(\alph*)]
  \setcounter{enumi}{6}
  \item The set of points whose distance to $a$ does not exceed a fixed fraction $\theta$ of the distance to $b$, i.e. $$\{x \ | \ \| x - a \|_2 \leq \theta \| x - b \|_2\},$$ where $a, b \in \mathbb{R}^n, a \neq b$ and $0 \leq \theta \leq 1$.
\end{enumerate}

\paragraph{Solution:}

Note that 

\begin{align*}
  \{x \ | \ \| x - a \|_2 \leq \theta \| x - b \|_2\} &= \{x \ | \ \| x - a \|_2^2 \leq \theta^2 \| x - b \|_2^2\} \\
  &= \{x \ | \ x^Tx - 2a^Tx + a^Ta \leq \theta^2 (x^Tx - 2b^Tx + b^Tb)\} \\
  &= \{x \ | \ x^Tx - 2a^Tx + a^Ta \leq \theta^2 x^Tx - 2\theta^2 b^Tx + \theta^2 b^Tb\} \\
  &= \{x \ | \ (1 - \theta^2) x^Tx - 2(a - \theta^2 b)^Tx + (a^Ta - \theta^2 b^Tb) \leq 0\}
\end{align*}

When $\theta = 1$, the set is a halfspace, which is convex. When $\theta \in [0,1)$, the set is an ball, which is convex. Therefore, the set $\{x \ | \ \| x - a \|_2 \leq \theta \| x - b \|_2\}$ is convex. $\square$





\section{Let $f: \mathbb{R} \rightarrow \mathbb{R}$. Show that $f$ is non decreasing if and only if there is a 
convex differentiable function $\varphi$ such that $f(x) = \varphi'(x)$.}

\paragraph{Proof:}

$\Leftarrow$ Suppose that there is a convex function $\varphi$ differentiable on $\mathbb{R}$ such that $f(x) = \varphi'(x)$. \\

Let $x_1, x_2, x_3, x_4 \in \mathbb{R}$ such that $x_1 \leq x_2 \leq x_3 \leq x_4$. Then, since $\varphi$ is convex, we have that

$$ \cfrac{\varphi(x_2)-\varphi(x_1)}{x_2-x_1} \leq \cfrac{\varphi(x_3)-\varphi(x_2)}{x_3-x_2} \leq \cfrac{\varphi(x_4)-\varphi(x_3)}{x_4-x_3}$$

Taking the limits as $x_2 \rightarrow x_1^{+}$ and $x_3 \rightarrow x_4^{-}$, by the definition of the derivative, we have that

$$ f(x_1) = \varphi'(x_1) \leq \varphi'(x_4) = f(x_4)$$

Therefore, $f$ is non decreasing.\\

$\Rightarrow$ Let $f$ be non decreasing and $\varphi$ a differentiable function such that $f(x) = \varphi'(x)$, i.e. $\varphi(x) = \int_{-\infty}^x f(t) \ dt$. It is clear that $f$ is integrable in the Lebesgue sense, because $f$ is non decreasing.\\


Let $x_1, x_2, x_3 \in \mathbb{R}$ such that $x_1 < x_2 < x_3$. By the mean value theorem, there exists $c_1, c_2$ such that 

\begin{align*}
  \cfrac{\varphi(x_2) -\varphi(x_1)}{x_2-x_1} = \varphi'(c_1) = f(c_1)\\
  \cfrac{\varphi(x_3) -\varphi(x_2)}{x_3-x_2} = \varphi'(c_2) = f(c_2)
\end{align*}

where $x_1 < c_1 < x_2 < c_2 < x_3$. Since $f$ is non decreasing, then $f(c_1) \leq f(c_2)$, then 

$$ \cfrac{\varphi(x_2) -\varphi(x_1)}{x_2-x_1} \leq \cfrac{\varphi(x_3) -\varphi(x_2)}{x_3-x_2}$$

Therefore, $\varphi$ is convex. $\square$


\section{Let $f: C \rightarrow \mathbb{R}$ a twice Fréchet differentiable function on a convex set $C \subseteq \mathbb{R}^n$. The following statements are known to be equivalent: Ex. 4.11. Foundations of Optimization, Güler.}

\begin{enumerate}[label=(\alph*)]
  \item $f$ is convex.
  \item $f(y) \geq f(x) + \langle \nabla f(x), y - x \rangle$ for all $x, y \in C$.
  \item $\langle \nabla f(y) - \nabla f(x), y - x \rangle \geq 0$ for all $x, y \in C$.
  \item $Hf(x)$ is positive semidefinite for all $x \in C$.
\end{enumerate}

Give direct proofs of 
\begin{enumerate}[label=(\alph*)]
  \item $\text{(b)} \Rightarrow \text{(c)}$
  \item $\text{(c)} \Rightarrow \text{(b)}$
  \item $\text{(c)} \Rightarrow \text{(d)}$
  \item $\text{(d)} \Rightarrow \text{(c)}$
\end{enumerate}

\paragraph{Proof:}


\begin{enumerate}[label=(\alph*)]
  \item $\text{(b)} \Rightarrow \text{(c)}$:
\end{enumerate}

Starting from the condition (b),

$$f(y) \geq f(x) + \langle \nabla f(x), y - x \rangle$$

Subtracting $f(x)$ from both sides, we have

$$f(y) - f(x) \geq \langle \nabla f(x), y - x \rangle$$

Now since this holds for all $x, y \in C$, it also follows that 

$$f(x) - f(y) \geq \langle \nabla f(y), x - y \rangle$$

Adding the last two inequalities, we have

$$f(y)- f(x) + f(x) - f(y) \geq \langle \nabla f(x), y - x \rangle + \langle \nabla f(y), x - y \rangle$$

Therefore,

$$0 \geq \langle \nabla f(x), y - x \rangle + \langle \nabla f(y), x - y \rangle = \langle \nabla f(x) - \nabla f(y), y - x \rangle$$

Which implies 

$$\langle \nabla f(y) - \nabla f(x), y - x \rangle \geq 0 \ \square$$ 






\begin{enumerate}[label=(\alph*)]
  \setcounter{enumi}{1}
  \item $\text{(c)} \Rightarrow \text{(b)}$:
\end{enumerate}


Starting from the condition (c),

$$\langle \nabla f(y) - \nabla f(x), y - x \rangle \geq 0$$

Let $x_{\tau} := x + \tau (y - x)$, then

\begin{align*}
  f(y) &= f(x) + \int_0^1 \langle \nabla f(x_{\tau}), y - x \rangle \ d\tau\\
  &= f(x) + \langle \nabla f(x), y - x \rangle + \int_0^1 \langle \nabla f(x_{\tau}) - \nabla f(x), y - x \rangle \ d\tau\\
  &= f(x) + \langle \nabla f(x), y - x \rangle + \int_0^1 \frac{1}{\tau} \langle \nabla f(x_{\tau}) - \nabla f(x), x_{\tau} - x \rangle \ d\tau\\
  &\geq f(x) + \langle \nabla f(x), y - x \rangle \ \square
\end{align*}


\begin{enumerate}[label=(\alph*)]
  \setcounter{enumi}{2}
  \item $\text{(c)} \Rightarrow \text{(d)}$:
\end{enumerate}

Starting from the condition (c),

$$\langle \nabla f(y) - \nabla f(x), y - x \rangle \geq 0$$

Let $y:=x+th$, then 

$$\langle \nabla f(x+th) - \nabla f(x), th \rangle \geq 0$$

Also, for the Hessian matrix, we have that

$$\nabla f(x+th) = \nabla f(x) + t Hf(x)h + o(|t|)$$

or equivalently

$$\nabla f(x+th) - \nabla f(x) = t Hf(x)h + o(|t|)$$

Therefore,

$$\langle \nabla f(x+th) - \nabla f(x), th \rangle = \langle t Hf(x)h + o(|t|), th \rangle = t^2 h^T Hf(x)h + o(t^2)$$

Thus, if $Hf(x)<0$, then the last expression is negative for a small enough $t$, a contradiction. Therefore, $Hf(x)$ is positive semidefinite. $\square$

\begin{enumerate}[label=(\alph*)]
  \setcounter{enumi}{3}
  \item $\text{(d)} \Rightarrow \text{(c)}$:
\end{enumerate}

Note that 

$$f(y) = f(x) + \nabla f(x)^T(y-x) + \cfrac{1}{2} (y-x)^T Hf(x)(y-x) + o(\|y-x\|^2)$$

Since $Hf(x)$ is positive semidefinite, then

$$f(y) \geq f(x) + \nabla f(x)^T (y-x)$$

which is the condition (b), which was proved to be equivalent to (c). $\square$



\setcounter{section}{6}

\section{Consider the function $f(x) = \langle c, x \rangle - \displaystyle\sum_{j=1}^m \log(1 - \langle a_j, x \rangle) - \displaystyle\sum_{i=1}^n \log(1- x_i^2)$. Show that $f$ is convex.}

\paragraph{Proof:}

Let $g(x) = \langle c, x \rangle$, $h(x) = \displaystyle\sum_{j=1}^m \log(1 - \langle a_j, x \rangle)$ and $l(x) = \displaystyle\sum_{i=1}^n \log(1- x_i^2)$.\\

$g(x)$ is linear, and therefore convex.\\

The function $1 - \langle a_j, x \rangle$ is affine, and therefore concave. Since the logarithm function is concave, then $h(x)$ is concave as it is the sum of concave functions.\\

As a function of one variable, $\log(1-x_i^2)$ has second derivative $- \cfrac{2(x_i^2 + 1)}{(1-x_i^2)^2}$, since for all $x_i \in \mathbb{R}$ both the numerator and the denominator are positive, then $(\log(1-x_i^2))'' < 0$, i.e. $\log(1-x_i^2)$ is concave. Therefore, $l(x)$ is concave as it is the sum of concave functions.\\

Since $g(x)$ is convex, $h(x)$ is concave, thus $-h(x)$ is convex, $l(x)$ is concave, then $-l(x)$ is convex, and the sum of convex functions is convex, then $f(x) = g(x) - h(x) - l(x)$ is convex. $\square$



\end{document}